% =============================================================================
% Appendix A: Technical Details
% =============================================================================

\subsection{Synthetic Network Generation}

Synthetic networks are generated using tile-based templates. Each template defines a connectivity pattern within a single tile; tiles are then arrayed to create networks of the desired extent.

\begin{description}
\item[Trellis] A regular grid with diagonal connections, representing highly connected urban cores. Average degree $\approx 4$.

\item[Tree] A hierarchical branching structure where each node connects to at most 3 others. Represents suburban cul-de-sac development. Average degree $\approx 2.3$.

\item[Linear] A constrained grid with fewer cross-connections than trellis. Intermediate connectivity. Average degree $\approx 3$.
\end{description}

Networks are tiled 24$\times$24 to achieve approximately 12km extent, then the largest connected component is extracted. Node positions are perturbed slightly to avoid perfectly regular spacing.

\subsection{GLA Network Preprocessing}

The Greater London Authority road network is sourced from Ordnance Survey Open Roads data (GeoPackage format).

Preprocessing steps:
\begin{enumerate}
\item Remove invalid and empty geometries
\item Explode multi-part geometries
\item Build NetworkX graph from edges
\item Remove filler nodes (degree-2 nodes that don't represent junctions)
\item Remove dangling nodes (dead-end segments $< 20$m)
\item Consolidate nearby nodes (threshold 1m)
\item Extract largest connected component
\end{enumerate}

Final network: $\approx$ \glaNnodes{} nodes.

\subsection{Implementation in \cityseer{}}
\label{sec:appendix_cityseer_impl}

The sampling model is implemented in the \cityseer{} Python library (version 4.23+). The adaptive function automatically probes reachability, computes sampling probabilities from the fitted model, and scales results:

\begin{verbatim}
from cityseer.metrics import networks

gdf = networks.node_centrality_shortest_adaptive(
    network_structure,
    nodes_gdf,
    distances=[500, 1000, 2000, 5000, 10000],
    compute_closeness=True,
    compute_betweenness=True,
    target_rho=0.95
)
\end{verbatim}

For manual control, a fixed sampling probability can be specified via the \texttt{sample\_probability} parameter of \texttt{node\_centrality\_shortest}. Helper functions for custom workflows (\texttt{compute\_required\_p}, \texttt{probe\_reachability}, \texttt{spatial\_sample}) are available in the \texttt{cityseer.config} module; see the library documentation for details.

\subsection{Reproducibility}

All analysis scripts, fitting data, and validation results are available at \texttt{https://github.com/benchmark-urbanism/cityseer-api/tree/master/analysis/sampling}. Model parameters are stored in \texttt{output/sampling\_model.json} and synced to \texttt{pysrc/cityseer/config.py} for library integration.
